\section{State Space, Determinism, and Noise}

\begin{definition}[State-space description]
A \textbf{state-space description} represents the system by a state vector
$x(t)$ together with an update rule. In continuous time this is written
\[
  \dot{x} = f(x,u,\eta),
\]
where $x$ denotes the current state, $u$ denotes control or intervention
inputs, and $\eta$ denotes stochastic or unmodeled influences.
\end{definition}

\begin{remark}
This equation is not offered as a fitted human model. It is a disciplined
template. It tells the reader what kinds of variables matter: current state,
external input, and noise. That is already a major improvement over vague talk
of ``feeling stuck.''
\end{remark}

The notation also helps integrate the earlier chapters. Chapter~\ref{ch:consciousness-biology}
supplied strong reasons to talk about organized state transitions in biology.
Chapter~\ref{ch:emergence} showed that repeated local interactions can form a
history-sensitive regime. The present chapter adds the mathematical language
needed to say what sort of regime that might be: a fixed point, a metastable
state, a limit cycle, a noisy transition path, or a control-sensitive branch of
system behaviour.
